\section*{Chapitre \ref{ch:mx:order-types-and-their-uses} --- Types d'ordres et leurs usages}

\begin{solution}{exo:mx:order-types-and-their-uses:1}
Il s'agit d'un \omterm{def:mx:order-types-and-their-uses:marketable}{ordre à cours limité exécutable}~: il achète 300 à 100.01 et 300 à 100.02, laissant 200 au cours vendeur de 100.02~; rien de l'ordre ne reste au repos.
\end{solution}

\begin{solution}{exo:mx:order-types-and-their-uses:2}
L'\omterm{def:mx:order-types-and-their-uses:tif}{ordre tout-ou-rien} demande 400 et n'en trouve que 300~: il est accepté puis annulé dans son intégralité, et rien ne s'échange. L'\omterm{def:mx:order-types-and-their-uses:tif}{ordre exécuter-ou-annuler} prend les 300
et annule les 100 restants.
\end{solution}

\begin{solution}{exo:mx:order-types-and-their-uses:3}
Les plateformes récompensent la liquidité affichée~: à un même prix, les ordres affichés passent devant les ordres non affichés quelle que soit leur heure d'arrivée. Un \omterm{def:mx:order-types-and-their-uses:hidden}{ordre caché} achète
le secret au prix de la priorité.
\end{solution}

\begin{solution}{exo:mx:order-types-and-their-uses:4}
$10\,000/500=20$ tranches apparaissent, et l'ordre se retrouve en fin de file 19 fois (chaque tranche après la première).
\end{solution}

\begin{solution}{exo:mx:order-types-and-their-uses:5}
$\sigma\sqrt T=0.5\sqrt{30}=2.74$ ticks, $d/\sigma\sqrt T=0.365$~: $C=2.74\,\varphi(0.365)-1\times(1-\Phi(0.365))=0.665$ ticks par action.
\end{solution}

\begin{solution}{exo:mx:order-types-and-their-uses:6}
Un \omterm{def:mx:order-types-and-their-uses:stop}{ordre stop à cours limité} qui atteint un carnet sans cours acheteur à son cours limite ou au-dessus se place en cours vendeur au lieu de vendre plus bas, si bien qu'il ne consomme pas les cours acheteurs suivants ni ne déclenche les
\omterm{def:mx:order-types-and-their-uses:stop}{ordres stop} suivants~: 25~cents au lieu de 37 avec vingt \omterm{def:mx:order-types-and-their-uses:stop}{ordres stop}. Son détenteur renonce à la certitude de sortir~: dans une vraie chute, l'ordre peut
rester non exécuté pendant que le cours continue de baisser.
\end{solution}

\begin{solution}{exo:mx:order-types-and-their-uses:7}
\texttt{iceberg\_study} avec \texttt{follow=True}~: le détecteur strict repère 43~\% des recharges au lieu de 62~\% (13 fausses alertes)~: une recharge au nouveau cours vendeur
n'est pas au prix de la tranche consommée. Le vendeur ne vend aussi que 23\,700 des 30\,000 actions sur la demi-heure, car les recharges qui suivent
un cours vendeur en baisse se placent derrière une nouvelle file au lieu de rencontrer des acheteurs à un prix que le marché a laissé.
\end{solution}

\begin{solution}{exo:mx:order-types-and-their-uses:8}
Un \omterm{def:mx:order-types-and-their-uses:peg}{ordre indexé au mid} (midpoint peg) ne fait gagner la moitié de la fourchette que lorsqu'il s'exécute~; or il ne s'exécute que si quelqu'un croise au mid, et c'est alors que ce quelqu'un en sait davantage
(sélection adverse, chapitre~9) ou qu'une contrepartie patiente se trouve justement là. Les ordres qui doivent s'exécuter entièrement prennent parfois de la liquidité~; un algorithme qui n'utilise que des \omterm{def:mx:order-types-and-their-uses:peg}{ordres indexés}
a un risque d'exécution et négocie surtout contre des flux informés.
\end{solution}

\begin{solution}{pb:mx:order-types-and-their-uses:1}
\textbf{1.} Les deux se placent derrière les ordres affichés à leur prix pour leur partie cachée~; l'\omterm{def:mx:order-types-and-their-uses:hidden}{ordre iceberg} affiche une tranche dont les recharges apparaissent sur le flux, alors que
l'\omterm{def:mx:order-types-and-their-uses:hidden}{ordre caché} n'affiche rien.
\textbf{2.} Une recharge est une augmentation de la taille affichée~: elle va en fin de file avec une nouvelle référence, comme toute nouvelle arrivée.
\textbf{3.} Il découle d'une exécution partielle qui laisse un reliquat non souhaité ou une position partiellement couverte~; il coûte les exécutions qu'il abandonne lorsque le carnet est peu fourni.
\textbf{4.} Chaque \omterm{def:mx:order-types-and-their-uses:stop}{ordre stop} déclenché devient un ordre de vente au marché qui consomme les cours acheteurs et imprime des cours plus bas, lesquels déclenchent les \omterm{def:mx:order-types-and-their-uses:stop}{ordres stop} suivants.
\textbf{5.} 5~cents~; 37~cents avec vingt \omterm{def:mx:order-types-and-their-uses:stop}{ordres stop}.
\textbf{6.} 19\,000 actions, pour une vente initiale de 3\,000.
\textbf{7.} 25~cents.
\textbf{8.} Davantage de profondeur entre les déclencheurs (fournisseurs de liquidité qui interviennent), un espacement plus large des \omterm{def:mx:order-types-and-their-uses:stop}{ordres stop}, ou des \omterm{def:mx:order-types-and-their-uses:stop}{ordres stop à cours limité}.
\textbf{9.} Après qu'une exécution retire entièrement un ordre affiché, signaler un ajout de même côté et au même prix dans les 50~ms dont la taille est égale à celle
initiale de l'ordre retiré.
\textbf{10.} 99 recharges, toutes détectées, avec 20 fausses alertes.
\textbf{11.} Sa recharge est un nouvel ordre qui gagne la plateforme et rencontre le marché tel qu'il est~: si le cours a bougé, une partie s'exécute dès son arrivée
et le reste se place plus petit ou à une autre position dans la file~; la recharge de la plateforme apparaît dans le même événement, à la même taille.
\textbf{12.} La détection tombe à 18~\% avec le délai et à 6~\% lorsque la taille est elle aussi rendue aléatoire.
\textbf{13.} 47~\% des recharges (34 détections), avec 378 fausses alertes~: onze fausses alertes pour chaque vraie.
\textbf{14.} \emph{Résultat récapitulatif}~: le détecteur strict repère 100~\% des recharges de plateforme, 62~\% des recharges algorithmiques immédiates, 18~\% avec un délai aléatoire de 200~ms
et 6~\% avec des tailles rendues aléatoires à hauteur de 30~\%~; le détecteur souple retrouve 47~\% de ces dernières au prix de 378 fausses alertes en une demi-heure.
\textbf{15.} Les tailles des tranches consommées au fil du temps (la même taille qui se répète), la persistance d'un \omterm{def:mx:the-limit-order-book:tick}{niveau de prix} et le déséquilibre d'exécutions à ce niveau.
\textbf{16.} Le temps~: un délai aléatoire prive le cours vendeur de sa liquidité, et les tranches se vendent plus tard~; sur un marché mouvant, certaines ne sont jamais exécutées.
\textbf{17.} Déduire l'intérêt caché à partir de données publiques est légal et courant~; négocier sur cette base relève de la concurrence ordinaire. Ce qui n'est pas légal, c'est de l'obtenir
auprès de la plateforme ou du courtier (le livre~9, chapitre~29, pour le bilan des sanctions en matière de manipulation~; le livre~16 pour les barrières à l'information).
\textbf{18.} Un \omterm{def:mx:order-types-and-their-uses:peg}{ordre indexé au mid} dans un dark pool, ou un algorithme qui fait varier la taille des tranches, le prix et le calendrier.
\textbf{19.} Rendre aléatoire la taille rechargée dans une plage donnée et retarder la recharge d'un délai aléatoire~; certaines plateformes proposent précisément ces options.
\textbf{20.} Son urgence, sa taille et sa patience~: chaque attribut est une information pour quiconque lit le flux.
\end{solution}

\begin{solution}{iq:mx:order-types-and-their-uses:1}
Lorsque la transaction doit se faire maintenant ou pas du tout~: couverture, prise d'un prix qui ne tiendra pas, balayage de plusieurs plateformes, sondage d'un dark pool. Un \omterm{def:mx:the-limit-order-book:orders}{ordre à cours limité}
au repos révèle son intérêt et risque de se faire cueillir par des contreparties mieux informées. \iqlookfor{le lien entre le type d'ordre, la fuite d'information et la sélection adverse.}
\end{solution}

\begin{solution}{iq:mx:order-types-and-their-uses:2}
Un ordre qui devient un \omterm{def:mx:the-limit-order-book:orders}{ordre au marché} lorsqu'une transaction s'imprime au prix stop ou au-delà~; il s'exécute alors contre tout ce que contient le carnet, ce qui, dans une
chute rapide, peut se situer bien plus bas (et les autres \omterm{def:mx:order-types-and-their-uses:stop}{ordres stop} s'ajoutent aux ventes). \iqlookfor{la conversion en \omterm{def:mx:the-limit-order-book:orders}{ordre au marché}, la cascade, l'alternative de l'\omterm{def:mx:order-types-and-their-uses:stop}{ordre stop à cours limité}.}
\end{solution}

\begin{solution}{iq:mx:order-types-and-their-uses:3}
Un ordre de vente au repos à $p$ est une option d'achat de prix d'exercice $p$ émise sur le marché~: celui qui apprend le premier que la valeur dépasse $p$ achète dessus. Un ordre d'achat au repos est un put.
L'option expire lorsque son détenteur l'annule, si bien que sa valeur croît avec le temps pendant lequel l'ordre reste sans surveillance et avec la volatilité. \iqlookfor{Copeland et
Galai~; pourquoi la rapidité d'annulation a une valeur monétaire.}
\end{solution}

\begin{solution}{iq:mx:order-types-and-their-uses:4}
Conserver la tranche affichée dans le carnet et la réserve dans l'enregistrement de l'ordre~; lorsque la tranche est entièrement exécutée et qu'il reste de la réserve, reporter la tranche suivante en
fin de niveau avec un nouvel horodatage~; sur le flux, une exécution de l'ancienne référence et un ajout sous une nouvelle (ou un remplacement). Décider et
documenter la priorité de la réserve par rapport aux autres ordres. \iqlookfor{priorité de la tranche rechargée, et le fait que le flux ne doit pas révéler la réserve.}
\end{solution}

\begin{solution}{iq:mx:order-types-and-their-uses:5}
La règle de Nasdaq, par exemple~: retirer l'ordre du carnet lorsque la fourchette interne est croisée ou absente, et le réinscrire au nouveau mid dès qu'un mid valide existe,
l'annuler si aucun ne se présente dans la seconde. Un mid défini par une cotation bloquée ou croisée n'est pas un prix équitable. \iqlookfor{que les prix indexés dépendent
d'une référence qui peut être invalide, et que le \omterm{def:mx:the-limit-order-book:engine}{moteur d'appariement} doit préciser ce qui se passe alors.}
\end{solution}

\begin{solution}{iq:mx:order-types-and-their-uses:6}
Chercher des schémas de recharge (la même taille réapparaissant au même prix après des exécutions), un niveau qui absorbe bien plus qu'il n'affiche, des
exécutions persistantes et unilatérales à un même prix, et des exécutions cachées (impressions de transactions non affichées) qui se regroupent à un prix~; pondérer chaque signal par son taux de base de
coïncidence. \iqlookfor{un détecteur concret et une prise en compte des fausses alertes.}
\end{solution}
